\SourceLeaf{191}
\par\Needspace{6\baselineskip}\medskip
\begin{center}\hypertarget{native-191-article142}{}\pdfbookmark[2]{142.}{article-142}142.\end{center}
Aequatio IV, quae evoluta ad ordinem octavum ascenderet, in forma sua non mutata expeditissime tentando solvitur. Ceterum e theoria aequationum facile ostendi potest (quod tamen fusius evolvere brevitatis caussa hic supersedemus), hanc aequationem vel duas vel quatuor solutiones per valores reales admittere\HistoricalNoteEight{Handschriftliche Bemerkung:] $z$ saltem duos valores reales habet, quoniam valores ipsius $G\sin z^4-\sin(z+H)$ pro $z=0$ atque pro $z=180^0$ signa opposita habent; $z$ non habet plures valores reales quam 4, quoniam $\dfrac{\sin(z+A)}{\sin z^4}$ inter $z=0$ atque $z=180^0$ semel tantum fit maximum semelque minimum, ac perinde inter $z=180^0$ et $360^0$. Scilicet hoc evenit, quoties $2\operatorname{tang}z=-3\operatorname{cotang}A\pm\BareRoot{(9\operatorname{cotang}A^2-16)}=8\operatorname{tang}(z+A)$. —
\par\indent Die Auflösung der Gleichung $a\sin z^4=\sin(z-A)$ lässt sich auf eine zierliche Construction zurückführen. Man setze $\operatorname{cotang}z=x$, also $\sin z=\dfrac{1}{\BareRoot{(1+xx)}}\mathpunct{\cdot}$ Folglich die Gleichung
\[
\frac{1}{(1+xx)^{\frac{3}{2}}}=\frac{\cos A}{a}-\frac{\sin A}{a}x.
\]
Es handelt sich also nun um den Schnitt der Curve $y=\dfrac{1}{(1+xx)^{\frac{3}{2}}}$ mit der Geraden $y=\dfrac{\cos A}{a}-\dfrac{\sin A}{a}x$.
\par\noindent\textsc{Binet}, \textit{Ecole polyt. Cah. 20. p. 285.}}. In casu priori valor alter ipsius $\sin z$ positivus erit, alterum negativum reiicere oportebit, quia per problematis naturam $r'$ negativus evadere nequit. In casu posteriori inter valores ipsius $\sin z$ vel unus positivus erit tresque reliqui negativi — ubi igitur haud ambiguum erit, quemnam adoptare oporteat — vel tres positivi cum uno negativo; in hoc casu e valoribus positivis ii quoque si qui adsunt reiici debent, ubi $z$ maior evadit quam $\delta'$, quoniam per aliam problematis conditionem essentialem $\rho'$ adeoque etiam $\sin(\delta'-z)$ quantitas positiva esse debet.

Quoties observationes mediocribus temporum intervallis ab invicem distant, plerumque casus postremus locum habebit, ut tres valores positivi ipsius $\sin z$ aequationi satisfaciant. Inter has solutiones praeter veram reperiri solet aliqua, ubi $z$ parum differt a $\delta'$, modo excessu, modo defectu: hoc phaenomenon sequenti modo explicandum est. Problematis nostri tractatio analytica ei soli conditioni superstructa est, quod tres corporis coelestis in spatio loci iacere debent \textit{in} rectis, quarum situs per locum absolutum terrae positionemque \SourceLeaf{192}observatam determinatur. Iam per ipsius rei naturam loci illi iacere quidem debent in \textit{iis} rectarum partibus, \textit{unde} lumen ad terram descendit: sed aequationes analyticae hanc restrictionem non agnoscunt, omniaque locorum systemata, qui quidem cum \textsc{Kepler}i legibus consentiunt, perinde complecti debent, sive ab hac terrae parte in illis rectis iaceant, sive ab illa, sive denique cum ipsa terra coincidant. Iam hic ultimus casus utique problemati nostro satisfaciet, quum terra ipsa ad normam illarum legum moveatur. Hinc patet, aequationes comprehendere debere solutionem, in qua puncta $C,C',C''$ cum punctis $A,A',A''$ coincidant (quatenus variationes minutissimas locis terrae ellipticis a perturbationibus et parallaxi inductas negligimus): aequatio itaque IV semper admittere deberet solutionem $z=\delta'$, si pro $P$ et $Q$ valores veri locis terrae respondentes acciperentur. Quatenus autem illis quantitatibus valores tribuuntur ab his non multum discrepantes (quod semper supponere licet, quoties temporum intervalla modica sunt), inter solutiones aequationis IV necessario aliqua reperiri debet, quae proxime ad valorem $z=\delta'$ accedit.

Plerumque quidem in eo casu, ubi aequatio IV tres solutiones per valores positivos ipsius $\sin z$ admittit, tertia ex his (praeter veram eamque de qua modo diximus) valorem ipsius $z$ maiorem quam $\delta'$ sistet, adeoque analytice tantum possibilis, physice vero impossibilis erit: tunc itaque quamnam adoptare oporteat ambiguum esse nequit. Attamen contingere utique potest, ut aequatio illa duas solutiones idoneas diversas admittat, adeoque problemati nostro per duas orbitas prorsus diversas satisfacere liceat. Ceterum in tali casu orbita vera a falsa facile dignoscetur, quamprimum observationes alias magis remotas ad examen revocare licuerit.
